\documentclass[11pt]{article}
\usepackage[margin=1in]{geometry}
\usepackage{booktabs}
\usepackage{hyperref}
\usepackage{amsmath,amssymb}
\usepackage{listings}
\usepackage{xcolor}
\lstset{basicstyle=\ttfamily\small,breaklines=true,frame=single,columns=fullflexible}

\title{QC-100 Independent Replication Report:\\
Spoofing Linear Cross-Entropy Benchmarking in Shallow Quantum Circuits\\
(arXiv:2005.02421)}
\author{Rick Stevens AI / OpenClaw replication pipeline (Ollie)}
\date{2026-07-03 (LaTeX conversion 2026-07-06)}

\begin{document}
\maketitle

\begin{abstract}
We independently reproduce the mechanistic core of Barak, Chou and Gao's
result that Google's linear cross-entropy benchmark (Linear XEB) is
classically spoofable at shallow depths, using CPU statevector
simulation on 1D brick-wall Haar-random circuits of $n\le 8$ qubits.
Uniform sampling yields $F_{\mathrm{XEB}}\approx 0$; exact sampling from
the ideal $q_C$ yields large positive $F_{\mathrm{XEB}}$ that decays
toward the Porter-Thomas limit $\to 1$ with depth. A trivial
\emph{light-cone-$L$} classical spoofer -- which never samples $q_C$ --
achieves $F_{\mathrm{XEB}}\approx 14$ at $n{=}8, d{=}1$ (light-cone
$L{=}2$), $\approx 2.3$ at $d{=}2$ (light-cone-4 block sampler), and
degrades smoothly as depth grows past the block budget, consistent with
the $1/15^d$ scaling of Theorem~1.1. Verdict: \textbf{REPLICATED} in the
small-instance regime. The full $\sqrt{\log n}$-depth 2D construction
(Corollary~1.2) and the 53-qubit Google-Sycamore number
$(2.24\pm 0.21)\cdot 10^{-3}$ are out of scope.
\end{abstract}

\section{Paper summary}

The Linear-XEB fidelity of a sampler $p$ on an $n$-qubit circuit $C$ is
\[
F_C(p)\;=\;\sum_x p(x)\bigl(2^n\,q_C(x) - 1\bigr),
\]
where $q_C(x) = |\langle x | \psi_C\rangle|^2$ is the ideal output
distribution. Exact sampling gives $F\approx 1$ for scrambled circuits;
uniform sampling gives $F\approx 0$; Google's 53-qubit device reported
$F\approx 2.24\cdot 10^{-3}$.

Barak, Chou and Gao prove (Theorem~1.1) that for a circuit of depth $d$
with Haar-random 2-qubit gates and per-output light-cone size $\le L$
there is a classical randomized algorithm running in $\mathrm{poly}(n,
2^L)$ time achieving
\[
\mathbb{E}\bigl[F_C(A_C)\bigr] \;=\; \Omega\bigl((L/n)\cdot 15^{-d}\bigr).
\]
Corollary~1.2 upgrades this to $F=\omega(1)$ for 2D circuits of depth
$d=O(\sqrt{\log n})$: the shallow-depth regime of Linear XEB is
\emph{classically spoofable}.

The paper is analytic and reports no numerical $F_{\mathrm{XEB}}$
values. The reproducible content is therefore the \emph{mechanism}:
(i) baseline behaviour of exact vs uniform sampling, (ii) Porter-Thomas
trajectory of the collision probability $\mathrm{CP}(q_C) = \sum_x
q_C(x)^2$ with depth, and (iii) a light-cone spoofer that beats the
$F\approx 0$ trivial-sampler line by orders of magnitude on shallow
circuits.

\section{Claims table}

\begin{center}\small
\begin{tabular}{clcc}
\toprule
ID & Claim & CPU-testable? & Tested here? \\
\midrule
C1 & Exact sampling $\Rightarrow F_{\mathrm{XEB}}\gg 0$ on shallow circuits & Yes ($n\le 8$) & Yes \\
C2 & Uniform sampling $\Rightarrow F_{\mathrm{XEB}}\approx 0$ & Yes & Yes \\
C3 & Light-cone-$L$ classical algorithm achieves $F_{\mathrm{XEB}}=\Omega(1)$ & Partial (mechanistic) & Yes \\
C4 & Advantage degrades with depth ($1/15^d$ scaling) & Yes & Yes \\
C5 & Full $\sqrt{\log n}$-depth 2D $\mathrm{poly}(n,2^L)$ construction & No (HPC-scale) & No \\
C6 & 53-qubit Sycamore $F=(2.24\pm 0.21)\cdot 10^{-3}$ (cited, not re-derived) & No (HPC-scale) & No \\
\bottomrule
\end{tabular}
\end{center}

\section{Method}

All experiments use a 1D brick-wall random circuit on $n$ qubits:
Haar-random $4\times 4$ unitaries on pairs $(0,1),(2,3),\ldots$ at even
depths and $(1,2),(3,4),\ldots$ at odd depths. Haar sampling follows
Mezzadri~(2007). The circuit is simulated with
\texttt{cirq.Simulator(dtype=complex128)}; $q_C(x) = |\langle
x|\psi_C\rangle|^2$; $\hat F = \frac{1}{N}\sum_{i=1}^N (2^n q_C(x_i) -
1)$ for $N$ i.i.d. samples $x_i\sim p$. Seed \texttt{20260703}, bit-for-bit
reproducible.

\paragraph{Environment.}
\begin{lstlisting}
python3 -m venv .venv
.venv/bin/pip install cirq-core numpy   # cirq-core 1.7.0, numpy 2.5.0
.venv/bin/python scripts/xeb_experiment.py               # ~4.5 min
.venv/bin/python scripts/collision_probability_check.py  # ~1 min
\end{lstlisting}

\paragraph{Spoofers.}
\begin{itemize}
\item \textbf{Depth-1 deterministic spoofer.} The brick-wall factorises
  into disjoint 2-qubit blocks $(0,1),(2,3),\ldots$; for each block
  compute the four marginals $q_k(v)$, pick $v^\star = \arg\max_v q_k(v)$,
  concatenate. Emit the same $x^\star$ every sample. This is the
  paper's Theorem~1.1 mechanism at $L{=}2$.
\item \textbf{Light-cone-4 block sampler ($d=2,3,6$).} Partition qubits
  into consecutive size-4 blocks, compute the exact marginal of $q_C$
  restricted to each block (by summing amplitudes), sample per block,
  stitch. At $d=2$ (true light-cone $\le 4$) this is the paper's
  strategy; at $d=3,6$ it is intentionally under-powered as a control.
\end{itemize}

\section{Results vs.\ paper}

\subsection{Exact and uniform baselines (C1, C2)}
$N=5000$ samples, 20 fresh random circuits per row.

\begin{center}\small
\begin{tabular}{ccrr}
\toprule
$n$ & $d$ & $F_{\mathrm{XEB}}$ exact & $F_{\mathrm{XEB}}$ uniform \\
\midrule
4 & 1 & $+1.383\pm 0.889$ & $-0.002\pm 0.011$ \\
4 & 6 & $+1.018\pm 0.495$ & $-0.005\pm 0.017$ \\
6 & 1 & $+2.707\pm 1.308$ & $-0.001\pm 0.020$ \\
6 & 6 & $+1.077\pm 0.295$ & $-0.005\pm 0.014$ \\
8 & 1 & $+5.193\pm 2.730$ & $-0.008\pm 0.038$ \\
8 & 6 & $+1.614\pm 0.524$ & $-0.001\pm 0.014$ \\
\bottomrule
\end{tabular}
\end{center}

Both C1 and C2 are confirmed quantitatively. Uniform $F_{\mathrm{XEB}}$
is $0\pm O(N^{-1/2})$; exact $F_{\mathrm{XEB}}$ is large and positive at
every depth and \emph{above} the Porter-Thomas value $\approx 1$ in the
shallow regime.

\subsection{Collision-probability trajectory (C4)}
$\mathbb{E}[F_{\mathrm{XEB}}(\text{exact})] = 2^n\cdot\mathrm{CP}(q_C) - 1$;
Porter-Thomas: $F_{\mathrm{PT}} = 2\cdot 2^n/(2^n+1) - 1$.

\begin{center}\small
\begin{tabular}{ccrrr}
\toprule
$n$ & $d$ & $\mathrm{CP}(q_C)$ & $2^n\cdot\mathrm{CP}-1$ & $F_{\mathrm{PT}}$ \\
\midrule
8 & 1  & 0.0303 & \textbf{6.75} & 0.9922 \\
8 & 2  & 0.0183 & 3.67          & 0.9922 \\
8 & 4  & 0.0131 & 2.35          & 0.9922 \\
8 & 6  & 0.0107 & 1.74          & 0.9922 \\
8 & 10 & 0.0084 & \textbf{1.16} & 0.9922 \\
6 & 1  & 0.0593 & \textbf{2.80} & 0.9692 \\
6 & 10 & 0.0316 & \textbf{1.02} & 0.9692 \\
\bottomrule
\end{tabular}
\end{center}

$\mathrm{CP}(q_C)$ decays smoothly toward $2/(2^n+1)$ with depth --
exactly the picture the paper draws in its introduction.

\subsection{Shallow-depth spoofer (C3, C4)}
$N=1000$, 20 circuits per row.

\begin{center}\small
\begin{tabular}{cclr}
\toprule
$n$ & $d$ & Spoofer & $F_{\mathrm{XEB}}$ \\
\midrule
6 & 1 & det.\ best-bitstring per 2-qubit block          & $\mathbf{+7.59\pm 3.88}$  \\
6 & 2 & light-cone-4 block sampler                       & $\mathbf{+2.24\pm 1.57}$  \\
6 & 3 & light-cone-4 (control, cone $>$ block)           & $+1.26\pm 0.72$           \\
6 & 6 & light-cone-4 (control)                           & $+0.66\pm 0.39$           \\
8 & 1 & det.\ best-bitstring per 2-qubit block          & $\mathbf{+13.88\pm 7.31}$ \\
8 & 2 & light-cone-4 block sampler                       & $\mathbf{+2.31\pm 1.22}$  \\
8 & 3 & light-cone-4 (control)                           & $+2.10\pm 1.31$           \\
8 & 6 & light-cone-4 (control)                           & $+0.77\pm 0.41$           \\
\bottomrule
\end{tabular}
\end{center}

This is the paper's headline mechanism, reproduced cleanly:
\begin{itemize}
\item At $d{=}1$ (light-cone $L{=}2$), a trivial deterministic spoofer
  achieves $F\approx 7.6$ ($n{=}6$) and $F\approx 13.9$ ($n{=}8$) --
  orders of magnitude above Google's reported $\approx 2\cdot 10^{-3}$.
\item At $d{=}2$, the light-cone-4 sampler clears $F\approx 2.3$.
\item At $d{=}3,6$, the fixed 4-qubit block strategy degrades smoothly,
  consistent with the $1/15^d$ scaling in Theorem~1.1.
\end{itemize}

\section{Verdict}

\textbf{REPLICATED} (small-instance CPU statevector regime).

The paper's central spoofing mechanism (C3, C4) is reproduced in a
faithful small-scale operationalisation, together with confirmation of
the exact-vs-uniform baselines (C1, C2) and the Porter-Thomas
trajectory. The full $\sqrt{\log n}$-depth 2D construction (C5) and the
53-qubit Sycamore number (C6) are HPC-scale and not attempted; the
paper itself does not verify either numerically.

\section{Honest critique of what we did (and did not) exercise}

\begin{itemize}
\item \textbf{The spoofing algorithm was independently reimplemented,
  but only in its simplest forms.} The depth-1 ``pick the argmax of
  each disjoint 2-qubit block'' spoofer is a full, honest reimplementation
  of Theorem~1.1 at $L{=}2$: no results borrowed from the paper are used
  as inputs. The depth-2 light-cone-4 block sampler is our own
  brute-force marginal computation (sum amplitudes over the complement),
  not the tensor-network reduction the paper introduces in \S3. Both
  match the paper's algorithmic \emph{scaling} ($\mathrm{poly}(n)\cdot
  2^L$) but the paper's construction is more sophisticated and would be
  the right target for a bigger reproduction.
\item \textbf{XEB scores against actual random circuits were reproduced,
  not quoted.} Every $F_{\mathrm{XEB}}$ number in \S4 is Monte-Carlo
  estimated from freshly generated Haar-random circuits with a fixed
  seed; nothing is taken from the paper (which reports none anyway).
\item \textbf{The shallow-depth threshold was verified only
  qualitatively.} We show the spoofer's advantage falling smoothly with
  depth for a fixed block size, which is consistent with the $1/15^d$
  factor in Theorem~1.1. We did \emph{not} fit the numerical decay to
  the $15^{-d}$ law, and we did not sweep enough $(n,d,L)$ points to
  test the $(L/n)$ prefactor. That would be a natural next probe.
\item \textbf{Comparison against genuine quantum sampling was not
  attempted.} The paper's numerical anchor is Google's $2.24\cdot
  10^{-3}$; matching that would require a 53-qubit noisy quantum device
  or an exact 53-qubit statevector simulation (HPC-scale, and the paper
  itself does not do it). Our comparison is against \emph{simulated}
  exact sampling from $q_C$, which is the correct classical baseline
  for the analytic claim.
\item \textbf{Ensemble mismatch.} The paper's Theorem~1.1 is
  architecture-agnostic given a light-cone bound; we use the 1D
  brick-wall of Corollary~1.2, not the 2D Sycamore-style lattice.
\item \textbf{$F>1$ is not a bug.} The Porter-Thomas ``$F\approx 1$''
  intuition is a large-$n$, deep-circuit limit. For shallow circuits
  $\mathrm{CP}(q_C)\gg 2/(2^n{+}1)$, so both exact sampling and any
  distribution-aware spoofer will exceed 1 -- exactly the loophole the
  paper exploits.
\end{itemize}

See also \texttt{report/failure\_analysis.md} and
\texttt{report/open\_questions.json} for a structured breakdown.

\section{Files}
\begin{lstlisting}
report/REPORT.md                         <- narrative markdown
report/REPORT.tex                        <- this LaTeX version
report/failure_analysis.md               <- honest critique
report/open_questions.json               <- 5 open questions
report/open_questions_section.tex        <- LaTeX version of questions
report/workflow.md                       <- reproducible workflow
report/artifacts_summary.md              <- file inventory
report/evidence/xeb_results.json         <- structured results
report/evidence/run.log                  <- console log
report/evidence/collision_check.log      <- CP trajectory log
scripts/xeb_experiment.py                <- main driver
scripts/collision_probability_check.py   <- CP sanity check
work/paper.pdf, work/paper.txt, work/abs.html
extraction/nougat.mmd                    <- extraction placeholder
.venv/                                   <- cirq-core 1.7.0, numpy 2.5.0
\end{lstlisting}

\end{document}
